% ---------------------------------------------------------------------
%  CHƯƠNG 3 — PHÂN TÍCH VÀ THIẾT KẾ
% ---------------------------------------------------------------------
\chapter{Phân tích và thiết kế}
\label{ch:thietke}

\begin{muctieu}
Chương này phân tích yêu cầu của bài toán, sau đó lần lượt trình bày thiết kế
kiến trúc hệ thống, thiết kế dữ liệu và đặc trưng, thiết kế các mô hình học,
thiết kế phần phục vụ và quy trình chọn mô hình.
\end{muctieu}

\section{Phân tích bài toán và yêu cầu}

\subsection{Đầu vào và đầu ra}

\textbf{Đầu vào} là 4 quan trắc gần nhất của một cơn bão, cách nhau 6 giờ,
tức 24 giờ lịch sử. Mỗi quan trắc gồm thời điểm UTC, vĩ độ, kinh độ, gió, khí
áp và cấp cường độ.

\textbf{Đầu ra} gồm ba phần: vị trí tâm bão tại +6, +12, +24, +48 và +72 giờ;
bán kính nón bất định tại từng mốc; và danh sách 5 cơn bão lịch sử tương tự
nhất, dùng làm lời giải thích cho gợi ý.

\subsection{Yêu cầu chức năng}

\begin{itemize}[leftmargin=1.6em, itemsep=2pt]
  \item Nạp và làm sạch dữ liệu best-track, tạo đặc trưng, huấn luyện mô hình.
  \item Dự báo quỹ đạo qua API cho bão đang hoạt động và cho chuỗi quan trắc
        do người dùng gửi lên.
  \item Hiển thị quỹ đạo dự báo, nón bất định và bão tương tự trên bản đồ.
\end{itemize}

\subsection{Yêu cầu phi chức năng}

\begin{itemize}[leftmargin=1.6em, itemsep=2pt]
  \item \textbf{Không rò rỉ dữ liệu} giữa huấn luyện và kiểm thử.
  \item \textbf{Tái lập được}: cố định seed, lưu cấu hình và mã băm dữ liệu
        của từng lần chạy.
  \item \textbf{Toàn vẹn}: mô hình được đóng gói kèm mã băm, phần phục vụ từ
        chối gói bị sửa đổi hoặc sai phiên bản.
  \item \textbf{Minh bạch}: mọi kết quả hiển thị kèm phiên bản mô hình, nguồn
        dữ liệu và cảnh báo ``chỉ mang tính tham khảo''.
\end{itemize}

\subsection{Tiêu chí chạy đúng}

Hệ thống được coi là đạt khi sai số Haversine trên tập kiểm thử thấp hơn
persistence có ý nghĩa thống kê, và nón 80\% phủ khoảng 80\% vị trí thực.

\section{Thiết kế kiến trúc hệ thống}

Kiến trúc tổng thể được trình bày ở Hình~\ref{fig:kientruc}. Hệ thống tách
thành hai nửa có vòng đời khác nhau. Nửa huấn luyện chạy ngoại tuyến, đi từ
dữ liệu thô đến gói mô hình. Nửa phục vụ chạy trực tuyến, nạp gói mô hình đó
để trả lời yêu cầu dự báo. Hai nửa chỉ giao tiếp qua gói mô hình, nên có thể
huấn luyện lại mà không sửa phần phục vụ.

\begin{figure}[H]
\centering
\resizebox{\textwidth}{!}{%
\begin{tikzpicture}[node distance=10mm and 9mm]
  \node[objbox] (data) {IBTrACS WP\\ 1980--2024};
  \node[clsbox, right=of data] (clean) {Làm sạch,\\ chuẩn hóa};
  \node[clsbox, right=of clean] (feat) {Tạo đặc trưng\\ (47 + 11 biến)};
  \node[clsbox, right=of feat] (split) {Cửa sổ chuỗi,\\ chia theo năm};
  \node[clsbox, right=of split] (train) {Huấn luyện,\\ chọn trên val};
  \node[objbox, right=of train] (bundle) {Gói mô hình\\ + thư viện analog};
  \node[objbox, below=of feat] (env) {NCEP/NCAR R1\\ (dòng dẫn đường)};
  \node[clsbox, below=of bundle] (api) {FastAPI\\ dự báo + nón};
  \node[objbox, left=of api] (web) {Bản đồ web\\ + bão tương tự};
  \draw[flowarr] (data) -- (clean);
  \draw[flowarr] (clean) -- (feat);
  \draw[flowarr] (feat) -- (split);
  \draw[flowarr] (split) -- (train);
  \draw[flowarr] (train) -- (bundle);
  \draw[flowarr] (env) -- (feat);
  \draw[flowarr] (bundle) -- (api);
  \draw[flowarr] (api) -- (web);
\end{tikzpicture}}
\caption{Sơ đồ kiến trúc tổng thể của hệ thống.}
\label{fig:kientruc}
\end{figure}

Mã nguồn chia thành các mô-đun độc lập, mỗi mô-đun ứng với một khối trong sơ
đồ:
\begin{itemize}[leftmargin=1.6em, itemsep=2pt]
  \item \code{ingestion}: tải và phân tích dữ liệu nguồn, ghi lại nguồn gốc.
  \item \code{features}: làm sạch, tạo đặc trưng, chuẩn hóa.
  \item \code{datasets}: cắt cửa sổ chuỗi và chia tập.
  \item \code{models}, \code{training}: định nghĩa mô hình, vòng huấn luyện,
        checkpoint và đóng gói.
  \item \code{evaluation}: tính sai số và bootstrap.
  \item \code{inference}: nạp gói mô hình, k-NN analog, nón bất định.
  \item \code{api}: giao diện HTTP và phục vụ trang bản đồ.
\end{itemize}

\section{Thiết kế dữ liệu}

\subsection{Nguồn và làm sạch}

Nguồn chính là NOAA IBTrACS v04r01~\cite{Knapp2010}, lưu vực Tây Bắc Thái Bình
Dương (WP), giai đoạn 1980--2024. Quy trình làm sạch gồm bốn bước: chuẩn hóa
đơn vị gió từ knot sang m/s; khử bản ghi trùng lặp; chỉ giữ quan trắc synop 6
giờ (00, 06, 12, 18 UTC); và giới hạn vùng 0--30°N, 100--140°E. Sau làm sạch
còn \textbf{24.024 quan trắc của 1.089 cơn bão}. Mỗi quan trắc lưu nguồn, URL
và mã băm SHA-256 của tệp gốc, nên có thể truy ngược bất kỳ con số nào về tệp
dữ liệu đã sinh ra nó.

\subsection{Chia tập không rò rỉ}

Các quan trắc liên tiếp của cùng một cơn bão rất giống nhau. Nếu chia ngẫu
nhiên theo quan trắc, mô hình sẽ gặp lại gần như cùng một tình huống ở cả
huấn luyện lẫn kiểm thử và cho điểm số đẹp giả tạo. Vì vậy dữ liệu được chia
\textbf{theo năm, giữ nguyên cả cơn bão} (Bảng~\ref{tab:chiatap}). Bộ chuẩn
hóa đặc trưng chỉ được fit trên tập huấn luyện rồi dùng lại nguyên vẹn cho
validation, test và khi suy luận.

\begin{table}[H]
\centering
\caption{Chia tập dữ liệu theo năm.}
\label{tab:chiatap}
\begin{tabularx}{\textwidth}{>{\bfseries\color{daunavy}}l >{\color{daumaroon}}c c X}
\toprule
\rowcolor{gray100}\bfseries\color{daunavy}Tập & \bfseries\color{daumaroon}Giai đoạn & \bfseries\color{daunavy}Số bão & \leavevmode\bfseries\color{daunavy}Mục đích \\
\midrule
Huấn luyện  & 1980--2016 & 895 & Học tham số ($\approx$7.800 cửa sổ) \\
Validation  & 2017--2020 & 110 & Chọn mô hình, siêu tham số, bán kính nón \\
Kiểm thử    & 2021--2024 & 84  & Đánh giá một lần (714 cửa sổ) \\
Ảnh hưởng VN & 2021--2024 & 27 & Tập con của kiểm thử: đi vào trong 300~km bờ biển (278 cửa sổ) \\
\bottomrule
\end{tabularx}
\end{table}

\subsection{Đặc trưng}

Một bộ tạo đặc trưng duy nhất (\code{FeatureBuilder}) dùng chung cho huấn
luyện và suy luận, nhờ đó hai giai đoạn không thể lệch nhau về cách tính. Mỗi
quan trắc có 47 đặc trưng, tóm tắt ở Bảng~\ref{tab:dactrung}. Thí nghiệm mở
rộng ở mục~\ref{sec:steering} bổ sung 11 đặc trưng dòng dẫn đường lấy từ tái
phân tích NCEP/NCAR R1~\cite{Kalnay1996}.

\begin{table}[H]
\centering
\caption{Các nhóm đặc trưng đầu vào.}
\label{tab:dactrung}
\begin{tabularx}{\textwidth}{>{\bfseries\color{daunavy}}l X}
\toprule
\rowcolor{gray100}\bfseries\color{daunavy}Nhóm & \leavevmode\bfseries\color{daunavy}Đặc trưng \\
\midrule
Vị trí     & Vĩ độ, kinh độ, chênh lệch và độ trễ 1--3 bước \\
Động học   & Quãng đường bước, hướng di chuyển, tốc độ (km/h), gia tốc, góc rẽ, tốc độ rẽ \\
Thời gian  & sin/cos của giờ, ngày, tháng, ngày trong năm; cờ mùa bão \\
Địa lý     & Khoảng cách, hướng tới bờ biển Việt Nam và các điểm tham chiếu (Hoàng Sa, Trường Sa, Đà Nẵng, \dots) \\
Cường độ   & Gió (m/s), khí áp (hPa), biến thiên của chúng, mã one-hot cấp bão \\
Dòng dẫn đường (mở rộng) & Dòng dẫn đường lớp sâu trong vành 200--800~km, gió 500~hPa, đứt gió 200--850~hPa, độ cao 500~hPa và chênh lệch 5° về bốn hướng \\
\bottomrule
\end{tabularx}
\end{table}

\section{Thiết kế mô hình}

\subsection{Biểu diễn đầu ra: độ dịch chuyển}

Thay vì để mạng dự báo trực tiếp tọa độ tuyệt đối, mô hình dự báo
\term{độ dịch chuyển} so với vị trí quan trắc cuối $\mathbf{p}_t$:
\begin{equation}
\hat{\mathbf{p}}_{t+h} = \mathbf{p}_t + \Delta\hat{\mathbf{p}}_h .
\end{equation}
Vị trí hiện tại $\mathbf{p}_t$ là thông tin đã biết chính xác. Nếu học tọa độ
tuyệt đối trên đầu vào đã chuẩn hóa, mạng phải tốn sức ``khôi phục'' lại chính
vị trí này. Học độ dịch chuyển cho phép mạng tập trung vào phần thực sự cần
dự báo là hướng và tốc độ di chuyển. Mục~\ref{sec:ablation} kiểm chứng lựa
chọn này bằng thí nghiệm.

\subsection{Các mô hình được so sánh}

Bảy phương pháp thuộc bốn nhóm được đặt cạnh nhau (Bảng~\ref{tab:phuongphap}).
Hai mốc so sánh không cần huấn luyện và cho biết mức sai số ``miễn phí''; mọi
phương pháp học phải vượt qua mốc này mới có giá trị.

\begin{table}[H]
\centering
\caption{Các phương pháp được huấn luyện và so sánh.}
\label{tab:phuongphap}
\begin{tabularx}{\textwidth}{>{\bfseries\color{daunavy}}l >{\color{daumaroon}}l X}
\toprule
\rowcolor{gray100}\bfseries\color{daunavy}Nhóm & \bfseries\color{daumaroon}Phương pháp & \leavevmode\bfseries\color{daunavy}Cấu hình \\
\midrule
Mốc so sánh & Persistence & Giữ nguyên vận tốc 6 giờ gần nhất \\
Mốc so sánh & CLIPER-proxy & Persistence có suy giảm \\
AI cổ điển & k-NN analog & $k = 40$; trọng số quỹ đạo 1,0, vị trí 1,0, mùa 0,5 \\
Học máy & Gradient boosting & HistGradientBoosting, hồi quy độ dịch chuyển \\
Học sâu & LSTM & 2 tầng, 128 nơ-ron ẩn, dropout 0,2 \\
Học sâu & Attention-LSTM & LSTM + attention cộng trên chuỗi đầu vào \\
Học sâu & Transformer & Encoder self-attention + mã hóa vị trí \\
\bottomrule
\end{tabularx}
\end{table}

Trong các phương pháp trên, k-NN analog có vai trò riêng: ngoài việc đưa ra
gợi ý, nó chỉ ra những cơn bão lịch sử làm căn cứ cho gợi ý đó.
Giải thuật~\ref{alg:knn} mô tả cách thực hiện.

\begin{algorithm}[H]
\caption{Gợi ý đường đi bằng k-NN analog}
\label{alg:knn}
\hspace*{\algorithmicindent}\textbf{Đầu vào:} cửa sổ quan trắc $q$; thư viện tình huống lịch sử $\mathcal{L}$; $k$; trọng số $w_g$\\
\hspace*{\algorithmicindent}\textbf{Đầu ra:} quỹ đạo gợi ý $\hat{\mathbf{y}}$ và danh sách bão tương tự
\begin{algorithmic}[1]
\Function{GoiYAnalog}{$q, \mathcal{L}, k$}
  \State $\mathbf{p}_t \gets$ vị trí cuối của $q$
  \State Tính đặc trưng $\mathbf{z}_g(q)$: đoạn quỹ đạo tương đối so với $\mathbf{p}_t$, vị trí, mùa
  \For{mỗi tình huống $c \in \mathcal{L}$}
    \State Tính $D(q,c)$ theo công thức~\eqref{eq:knn-dist}
  \EndFor
  \State $N \gets k$ tình huống có $D$ nhỏ nhất
  \State $\alpha_j \gets \dfrac{1/(D_j+\epsilon)}{\sum_{i \in N} 1/(D_i+\epsilon)}$ với mọi $j \in N$
  \State $\hat{\mathbf{y}} \gets \mathbf{p}_t + \sum_{j \in N} \alpha_j \, \Delta\mathbf{y}_j$
  \State \Return $\hat{\mathbf{y}}$, tên và năm của 5 cơn bão khác nhau gần $q$ nhất
\EndFunction
\end{algorithmic}
\end{algorithm}

\subsection{Cấu hình huấn luyện mạng nơ-ron}

\begin{itemize}[leftmargin=1.6em, itemsep=2pt]
  \item Hàm mất mát Smooth~L1 trên độ dịch chuyển (đơn vị độ).
  \item Adam, tốc độ học $10^{-3}$, weight decay $10^{-5}$, lịch cosine,
        cắt gradient ở chuẩn 1,0.
  \item Tối đa 100 epoch, dừng sớm khi validation không cải thiện sau 10 epoch,
        giữ trọng số của epoch tốt nhất trên validation.
  \item Mỗi mạng chạy 3 seed (0, 1, 2), báo cáo trung bình $\pm$ độ lệch chuẩn.
\end{itemize}

\section{Thiết kế phần phục vụ}

\subsection{Gói mô hình}

Kết quả huấn luyện được xuất thành một thư mục gồm trọng số, tham số của bộ
chuẩn hóa, bản kê chia tập và một tệp \code{manifest} ghi phiên bản, cấu hình
kiến trúc, danh sách đặc trưng, các mốc dự báo, bán kính nón và mã băm SHA-256
của từng tệp. Khi nạp, phần phục vụ kiểm tra lần lượt phiên bản lược đồ, mã
băm từng tệp, danh sách đặc trưng và kiểu mô hình thuộc danh sách cho phép;
sai một bước là từ chối nạp. Trọng số chỉ được đọc ở chế độ dữ liệu thuần,
không thực thi mã đi kèm tệp.

\subsection{Giao diện API}

Phần phục vụ là một ứng dụng FastAPI. Mỗi phản hồi dự báo mang theo phiên bản
mô hình, phiên bản dữ liệu, nguồn quan trắc, phương pháp tính bất định, các
cảnh báo (ví dụ quan trắc cũ hơn 24 giờ) và dòng tuyên bố miễn trừ. Các lỗi
trả về theo một khuôn dạng thống nhất gồm mã lỗi, thông điệp và mã yêu cầu để
tra nhật ký. Những đường dẫn ghi dữ liệu hoặc chạy mô hình yêu cầu khóa API
hoặc token, và các đường dẫn dữ liệu đều bị giới hạn tần suất theo địa chỉ
gọi.

\subsection{Nguyên tắc an toàn khi phục vụ}

\begin{itemize}[leftmargin=1.6em, itemsep=2pt]
  \item Ở chế độ vận hành, hệ thống từ chối khởi động nếu thiếu thông tin xác
        thực hoặc vẫn cho phép phương pháp dự phòng, thay vì lặng lẽ chạy
        không có mô hình đã duyệt.
  \item Đầu vào bị từ chối nếu quan trắc không theo thứ tự thời gian, có
        thời điểm ở tương lai hoặc cho tốc độ di chuyển phi thực tế.
  \item Bộ nhớ đệm chỉ là phần tăng tốc: khi nó hỏng, dự báo vẫn được tính
        lại từ đầu và kết quả không đổi.
\end{itemize}

\section{Quy trình chọn mô hình và nón bất định}

Mọi lựa chọn, gồm kiến trúc, đặc trưng, siêu tham số và thành viên của mô
hình kết hợp, chỉ dựa trên tập validation 2017--2020. Tập kiểm thử 2021--2024
được chạy \textbf{một lần} sau khi đã chốt cấu hình. Bán kính nón bất định
tại mỗi mốc lấy bằng phân vị 80\% của sai số trên validation.

\begin{luuy}
Gradient boosting có dòng dẫn đường đạt điểm test tốt nhất trong một thí
nghiệm, nhưng kém hơn trên validation nên không được chọn. Chọn mô hình theo
điểm test sẽ làm con số báo cáo lạc quan hơn thực tế.
\end{luuy}
